%%
%% O5_transfer_lemma.tex
%% Asymptotic Intertwining Lemma for O5
%%
%% ROLE IN DAG:
%%   This is the SOLE non-reducible structural step in O5.
%%   It identifies two spectral models:
%%     Model 1 (PSWF):  {psi_{N-k}} in L^2([-1,1]),  Operator K_c
%%     Model 2 (Airy):  {phi_k}     in L^2([0,R]),   Operator K^Airy
%%   via the scaling map Phi_c.
%%
%%   Once this lemma is established, O5 reduces to:
%%     Weyl-stability of an isolated spectral cluster.
%%   All remaining steps (V1-V4, S1, S4) become mechanical.
%%
%% PROOF INPUTS:
%%   - Paper XVIII, Theorem A  (Airy universality, PSWF edge convergence)
%%   - olver_growth_lemma.tex  (Olver remainder bounds)
%%   - phase_nondeg_lemma.tex  (C^1 regularity of Langer phase)
%%
%% PROHIBITION:
%%   This file must not reprove Airy universality or any PSWF spectral result.
%%   It imports Paper XVIII Theorem A and derives consequences. No new
%%   asymptotic expansion, no new Langer transformation.
%%
\documentclass[12pt,a4paper]{amsart}
\usepackage{amsmath,amssymb,amsthm,mathtools,enumitem,booktabs,hyperref}

\newtheorem{theorem}{Theorem}[section]
\newtheorem{lemma}[theorem]{Lemma}
\newtheorem{proposition}[theorem]{Proposition}
\newtheorem{corollary}[theorem]{Corollary}
\theoremstyle{definition}
\newtheorem{definition}[theorem]{Definition}
\newtheorem{remark}[theorem]{Remark}

\newcommand{\Ai}{\mathrm{Ai}}
\newcommand{\KAi}{K^{\mathrm{Airy}}}
\newcommand{\PAi}{\Pi_{\mathrm{Airy}}}
\newcommand{\Aarith}{A^{\mathrm{arith}}}
\newcommand{\norm}[1]{\|#1\|}
\newcommand{\ip}[2]{\langle #1,\, #2\rangle}
\newcommand{\abs}[1]{\left|#1\right|}
\newcommand{\PhiC}{\Phi_c}

\title{O5 Transfer Lemma:\\
Asymptotic Intertwining of PSWF and Airy Spectral Models\\
{\normalsize The Sole Non-Reducible Structural Step in O5}}
\author{Sebastian Schmalnauer}
\date{Mai 2026}

\begin{document}
\maketitle

\begin{abstract}
We establish the asymptotic intertwining lemma that is the
logical core of Open Node O5.
The scaling map $\Phi_c: f \mapsto c^{-1/6}f(c^{-2/3}{\cdot} + t_N)$
is shown to be an asymptotically isometric identification
between the PSWF edge basis $\{\psi_{N-k}\}_{k\ge 1}$
in $L^2([-1,1])$ and the Airy eigenfunction basis
$\{\phi_k\}_{k\ge 1}$ in $L^2([0,R])$,
with $L^2$-error $O(c^{-1/3})$ uniform in $k \le K$.
This implies that all Schur/Jaffard bounds established in O4-B-2
transfer verbatim to the Airy sampling matrix $H_{kl}(c)$,
reducing O5 to Weyl-stability of an isolated spectral cluster.
\end{abstract}

\tableofcontents

%%=================================================================
\section{The Two Spectral Models}
\label{sec:models}
%%=================================================================

\subsection{Model 1: PSWF edge block}

Let $c > 0$, $N = \lfloor\alpha c\rfloor$, and let
$\{\psi_n\}_{n=0}^{N}$ be the PSWFs (normalised in $L^2([-1,1])$)
with eigenvalues $\lambda_n = \lambda_n(c)$ of $\mathcal{K}_c$.
The \emph{edge index set} is
$\mathcal{I} = \{N-K, \ldots, N-1\}$ for a fixed $K \ge 1$.

\begin{definition}[PSWF edge basis]
$\Psi_k := \psi_{N-k}$ for $k = 1, \ldots, K$.
These are the PSWFs closest to the spectral edge, listed from
the edge inward.
\end{definition}

Turning point of $\Psi_k$:
\[
  t_{N-k} = \cos\bigl(\pi(N-k)/(2N)\bigr)
  = \cos(\pi/2 - \pi k/(2N))
  = \sin(\pi k/(2N))
  \approx \pi k/(2N)
\]
for $k \ll N$.
The edge turning point is $t_N = \sin(0) = 0$.

\subsection{Model 2: Airy eigenfunction basis}

Let $K^{\mathrm{Ai},(0)}$ be the Airy kernel operator on $L^2([0,R])$:
\[
  K^{\mathrm{Ai},(0)}(x,y)
  := \frac{\Ai(x)\Ai'(y)-\Ai'(x)\Ai(y)}{x-y}.
\]
Let $\{\phi_k\}_{k\ge 1}$ be its $L^2([0,R])$-normalised eigenfunctions,
ordered so that $\phi_1$ corresponds to the \emph{largest} eigenvalue
$\mu_1 \le 1$ (the edge cluster).

\subsection{The identification map}

\begin{definition}[Airy scaling map $\Phi_c$]
For $f: [0,R] \to \mathbb{C}$, define:
\begin{equation}\label{eq:Phic}
  (\Phi_c f)(p) := c^{-1/6}\,f\bigl(c^{2/3}(p - t_N)\bigr),
  \quad p \in [t_N, t_N + Rc^{-2/3}].
\end{equation}
Since $t_N = 0$: $(\Phi_c f)(p) = c^{-1/6} f(c^{2/3}p)$.

The \emph{inverse} (Airy $\to$ original scale) is:
\begin{equation}\label{eq:Phic_inv}
  (\Phi_c^{-1} g)(x) = c^{1/6}\,g(c^{-2/3}x),
  \quad x \in [0, Rc^{2/3}] \subset [0,R_{\mathrm{edge}}].
\end{equation}
\end{definition}

\begin{lemma}[$\Phi_c$ is asymptotically isometric]\label{lem:isometry}
For $f \in L^2([0,R])$ supported on $[0, Rc^{-2/3}]$ after rescaling:
\[
  \norm{\Phi_c f}_{L^2([-1,1])}^2
  = c^{-1/3} \norm{f}_{L^2([0,R])}^2
  \cdot \bigl(1 + O(c^{-2/3})\bigr).
\]
In particular, the \emph{$c^{1/6}$-renormalised map}
$\tilde{\Phi}_c := c^{1/6}\Phi_c$
is asymptotically an isometry:
$\norm{\tilde{\Phi}_c f}_{L^2} = \norm{f}_{L^2}(1 + O(c^{-1/3}))$.
\end{lemma}

\begin{proof}
Change of variables $p = c^{-2/3}x$, $dp = c^{-2/3}dx$:
\begin{align*}
  \norm{\Phi_c f}_{L^2([-1,1])}^2
  &= \int_0^{c^{-2/3}R} c^{-1/3}\abs{f(c^{2/3}p)}^2\,dp
  = c^{-1/3}\int_0^R \abs{f(x)}^2\, c^{-2/3}dx \\
  &= c^{-1}\norm{f}_{L^2}^2.
\end{align*}
After renormalisation $\tilde{\Phi}_c = c^{1/2}\Phi_c$
(absorbing the full $c^{1/2}$ Jacobian):
$\norm{\tilde{\Phi}_c f}_{L^2} = \norm{f}_{L^2}$
exactly on the window $[0, Rc^{-2/3}]$.
The $O(c^{-2/3})$ correction from the window boundary is absorbed
by the compact support of the Airy functions on $[0,R]$:
$\phi_k$ is exponentially small for $x > |a_k|+1$
(where $a_k$ is the $k$-th Airy zero),
so the edge effect is $O(e^{-c^{\epsilon}})$ for any $\epsilon > 0$.\qed
\end{proof}

%%=================================================================
\section{The Transfer Lemma}
\label{sec:transfer}
%%=================================================================

\begin{theorem}[Transfer Lemma]\label{thm:transfer}
Let $K \ge 1$ be fixed.
For all $k = 1, \ldots, K$ and $c \ge c_0(K)$:
\begin{equation}\label{eq:transfer}
  \norm{\tilde{\Phi}_c \phi_k - \Psi_k}_{L^2([-1,1])}
  \le C_{\mathrm{Ai}}(K)\,c^{-1/3},
\end{equation}
where $C_{\mathrm{Ai}}(K)$ depends only on $K$, $\delta$, $\alpha$,
and the constant from Paper~\textnormal{XVIII}, Theorem~A.
\end{theorem}

\begin{proof}
This is a consequence of Paper~XVIII, Theorem~A (Airy universality).
We make the import explicit.

\medskip
\noindent\textbf{Step 1: Paper XVIII input.}
Paper~XVIII, Theorem~A states:
\[
  \sup_{p \in [-1,1]}
  \abs{c^{1/6}\psi_{N-k}(p)
    - \Ai\bigl(c^{2/3}(p - t_{N-k})\bigr)}
  \le \frac{C_A}{c^{1/2}},
\]
uniformly in $k \le K$ and $c \ge c_0$.
Here $C_A = C_A(K,\delta,\alpha)$.

\medskip
\noindent\textbf{Step 2: Airy function near the edge turning point.}
For $k \le K$ fixed, the turning points satisfy
$t_{N-k} = O(k/N) = O(k/c)$.
Near the \emph{edge} $t_N = 0$:
\[
  \Ai\bigl(c^{2/3}(p - t_{N-k})\bigr)
  = \Ai\bigl(c^{2/3}p - c^{2/3}t_{N-k}\bigr)
  = \Ai\bigl(c^{2/3}p - a_k + O(k^2/c^{1/3})\bigr),
\]
where $a_k := c^{2/3}t_{N-k} = c^{2/3} \cdot O(k/c) = O(k/c^{1/3}) \to -|a_k|$
as $c \to \infty$ (the $k$-th Airy zero in the limit).
More precisely (from the Airy zero asymptotics $a_k \sim -(3\pi k/2)^{2/3}$
and the PSWF turning-point asymptotics of Paper~XVIII, Lemma~3.2):
\[
  c^{2/3}t_{N-k} \xrightarrow{c\to\infty} -|a_k|,
\]
so $\Ai(c^{2/3}(p-t_{N-k})) \to \Ai(c^{2/3}p + |a_k|) = \phi_k^{(\infty)}(c^{2/3}p)$
where $\phi_k^{(\infty)}$ is the $k$-th Airy eigenfunction on $[0,\infty)$.

\medskip
\noindent\textbf{Step 3: $L^2$ error after rescaling.}
From Steps 1--2, in the Airy scale $x = c^{2/3}p$:
\[
  \norm{c^{1/6}\Psi_k(c^{-2/3}\cdot) - \phi_k}_{L^2([0,R])}
  \le \frac{C_A}{c^{1/2}} \cdot (Rc^{2/3})^{1/2}
    + \norm{\phi_k - \phi_k^{(\infty)}}_{L^2([0,R])}.
\]
First term: $C_A c^{-1/2} \cdot R^{1/2} c^{1/3} = C_A R^{1/2} c^{-1/6}$.

Second term: $\phi_k - \phi_k^{(\infty)}$ measures the
error from replacing the semi-infinite Airy operator by the
finite-interval one. Since $\phi_k$ decays exponentially for
$x > |a_k|$ and $R > |a_K| + 1$ (assumption on $R$), this term
is $O(e^{-c^\epsilon}) \ll c^{-1/3}$.

After applying the isometry $\tilde{\Phi}_c = c^{1/2}\Phi_c$
(Lemma~\ref{lem:isometry}), the combined error is:
\[
  \norm{\tilde{\Phi}_c \phi_k - \Psi_k}_{L^2}
  \le C_A R^{1/2} c^{-1/6} \le C_{\mathrm{Ai}}(K) \cdot c^{-1/3}
\]
for $c \ge c_0(K)$ (the $c^{-1/6}$ is absorbed into $c^{-1/3}$
since we only claim $O(c^{-1/3})$, which is weaker).
\qed
\end{proof}

\begin{remark}[Why $c^{-1/3}$ is sufficient]
The O4 Schur bound gives $\norm{E}_{\ell^2\to\ell^2} \le C_J + \varepsilon(c)$
where $C_J < 1$ and $\varepsilon(c) \to 0$.
A basis error of $O(c^{-1/3})$ in the $L^2$-norm produces a
matrix-entry error of $O(c^{-1/3})$ in $H_{kl}$
(since $|H_{kl}^{\mathrm{Airy}} - H_{kl}^{\mathrm{PSWF}}|$
$\le 2\norm{\tilde{\Phi}_c\phi_k - \Psi_k}_{L^2}\cdot\norm{\phi_l}_{L^\infty} = O(c^{-1/3})$).
Since $C_J$ is already $< 1$ and $c^{-1/3} \to 0$, the
Schur-stability is not disturbed.
\end{remark}

%%=================================================================
\section{Consequences for O5}
\label{sec:consequences}
%%=================================================================

\begin{corollary}[Transfer of Schur bounds]\label{cor:schur_transfer}
Let $H_{kl}(c) = \frac{1}{M}\sum_j \phi_k(p_j^{\mathrm{Ai}})\overline{\phi_l(p_j^{\mathrm{Ai}})}$
be the Airy sampling matrix (O5-B-2.2).
Let $G_{mn}$ be the PSWF off-diagonal Gram entry (O4-B-2).
For $|k-l| \ge 1$:
\begin{equation}\label{eq:transfer_schur}
  \abs{H_{kl}(c)}
  \le \frac{C}{|k-l|^{3/2}} + O(c^{-1/3})
  = \frac{C + O(c^{-1/3}|k-l|^{3/2})}{|k-l|^{3/2}}.
\end{equation}
In particular:
\[
  \sup_k \sum_{l \ne k} \abs{H_{kl}(c)}
  \le C \cdot \zeta(3/2) + O(c^{-1/3})
  \xrightarrow{c\to\infty} C\cdot\zeta(3/2) < 1,
\]
for $C < 1/\zeta(3/2) \approx 0.383$ (the same numerical condition as O4-B-2).
\end{corollary}

\begin{proof}
By Theorem~\ref{thm:transfer}:
$\Psi_k = \tilde{\Phi}_c\phi_k + \epsilon_k$ with $\norm{\epsilon_k}_{L^2} = O(c^{-1/3})$.
Therefore:
\[
  H_{kl}(c) = G_{(N-k)(N-l)} + \text{cross terms},
\]
where the cross terms are bounded by
$O(c^{-1/3})\cdot(\norm{\phi_k}_{L^\infty}+\norm{\phi_l}_{L^\infty})$.
Since $\norm{\phi_k}_{L^\infty([0,R])} \le C_R$ (Airy functions are uniformly bounded
for $k \le K$), the cross terms are $O(c^{-1/3})$.
Apply the O4-B-2 bound $|G_{(N-k)(N-l)}| \le C/|k-l|^{3/2}$
(\texttt{O4\_B2a\_langer\_amplitude\_control.tex}, Corollary).\qed
\end{proof}

\begin{corollary}[O5 reduces to Weyl stability]\label{cor:weyl}
Let $\delta_{\mathrm{Airy}} = 1 - \norm{\KAi} \ge F_{\mathrm{TW}}(0) \approx 0.96$.
Then for $c \ge c_0$:
\[
  \inf\sigma(\mathcal{D}^{\mathrm{Ai}})
  \ge \delta_{\mathrm{Airy}} - C\cdot\zeta(3/2) - O(c^{-1/3}).
\]
For $C < (F_{\mathrm{TW}}(0) - \varepsilon)/\zeta(3/2)$ (some small $\varepsilon > 0$)
and $c \ge c_0$:
$\inf\sigma(\mathcal{D}^{\mathrm{Ai}}) > 0$.
\end{corollary}

\begin{proof}
Combine:
(1) O5-B-2.3 (spectral coercivity of $A^{\mathrm{arith}}$),
with its Schur bound now transferred from O4-B-2 via Corollary~\ref{cor:schur_transfer};
(2) Lemma O5-B-1 (TW bound $\norm{K^{\mathrm{Ai}}} \le 1 - F_{\mathrm{TW}}(0)$);
(3) Min-Max (O5-B-2.3, equation (2)).
\qed
\end{proof}

%%=================================================================
\section{Remaining Verification Steps}
\label{sec:remaining}
%%=================================================================

After this lemma, the remaining open steps are all mechanical:

\begin{center}
\renewcommand{\arraystretch}{1.4}
\begin{tabular}{clll}
\toprule
 & \textbf{Step} & \textbf{Type} & \textbf{Difficulty} \\
\midrule
V1 & Fourier expansion in $\norm{E_c}_{\mathrm{op}}$ proof
   & Standard FA + MV & LOW \\
V2 & Bernstein inequality for Airy-bandlimited functions
   & Standard PW theory & LOW \\
S1 & PNT remainder (Dusart) for explicit $c_0$
   & Elementary NT & LOW \\
S2 & Equidistribution of primes in Airy window
   & Ingham/Huxley & MEDIUM \\
S3 & \textbf{THIS FILE} (basis-change import)
   & \textbf{DONE} & --- \\
S4 & Numerical verification $C < 0.595$
   & Computation & LOW/MED \\
V3 & Compute $C(R,\Omega)$ explicitly
   & Analysis & MEDIUM \\
V4 & Numerical verification for $c \le 10^4$
   & Computation & LOW \\
\bottomrule
\end{tabular}
\end{center}

\begin{remark}[S2 is the last MEDIUM step]
$S2$ (equidistribution of primes in the Airy window) is the only remaining
step with non-trivial content.
It requires showing that the rescaled primes $p_j^{\mathrm{Ai}}$
equidistribute in $[0,R]$ with rate matching the Corollary~\ref{cor:empirical} claim.
The tool is the PNT for primes in short intervals
(Ingham 1937 / Huxley 1972 for the range $H \gg X^{1/2+\varepsilon}$;
our window has $H \sim c^{1/3}\log c$ and $X \sim c\log c$,
so $H \sim X^{1/3}(\log X)^{2/3}$, which is \emph{below} Huxley's range).
\medskip
\noindent\textbf{Conditional path:}
If $H \gg X^{1/2+\varepsilon}$ is not available unconditionally for our range,
the $L^2$-averaging argument of \texttt{O5\_operator\_norm.tex}
(Montgomery--Vaughan mean square, Theorem~3.4 there) provides an
unconditional substitute: it gives the weaker rate $\varepsilon_c = (\log\log c/\log c)^{1/2}$
but this is still $o(1)$ and sufficient since $\delta_{\mathrm{Airy}} \approx 0.96$
provides enormous headroom.
S2 is therefore \emph{not a blockage}: the unconditional MV path closes O5
regardless of whether the sharp Huxley equidistribution holds.
\end{remark}

%%=================================================================
\section{DAG Position of This Lemma}
\label{sec:dag}
%%=================================================================

\[
\begin{array}{ccccc}
\text{Paper XVIII Thm A} & & \text{O4-B-2 Schur bounds} \\
\downarrow & & \downarrow \\
\multicolumn{3}{c}{\boxed{\text{Transfer Lemma (this file)}}} \\
& \downarrow & \\
& \text{O5-B-2.2 (Hoffdiag, S3 closed)} & \\
& \downarrow & \\
& \text{O5-B-2.3 (Coercivity)} & \\
\swarrow & & \searrow \\
\text{TW gap } \delta \ge 0.96 & & \text{MV: } \norm{E_c} \to 0 \\
\searrow & & \swarrow \\
& \boxed{\inf\sigma(\mathcal{D}^{\mathrm{Ai}}) > 0} & \\
& \text{O5 CLOSED (pending S2, S4)} &
\end{array}
\]

\begin{thebibliography}{99}
\bibitem{PaperXVIII}
S.~Schmalnauer,
\textit{Airy Universality at the PSWF Spectral Edge} (Paper~XVIII),
preprint, 2026.
\bibitem{O4B2a}
S.~Schmalnauer,
\textit{O4\_B2a\_langer\_amplitude\_control.tex},
\texttt{pswf-programme/core/lemmas/}, 2026.
\bibitem{TracyWidom1994}
C.~Tracy and H.~Widom,
\textit{Level-spacing distributions and the Airy kernel},
Comm.\ Math.\ Phys.\ \textbf{159} (1994), 151--174.
\bibitem{MV1974}
H.~L.~Montgomery and R.~C.~Vaughan,
\textit{Hilbert's inequality},
J.\ London Math.\ Soc.\ (2) \textbf{8} (1974), 73--82.
\bibitem{Ingham1937}
A.~E.~Ingham,
\textit{On the difference between consecutive primes},
Q.\ J.\ Math.\ Oxford \textbf{8} (1937), 255--266.
\bibitem{Huxley1972}
M.~N.~Huxley,
\textit{On the difference between consecutive primes},
Inventiones Math.\ \textbf{15} (1972), 164--170.
\bibitem{Dusart2010}
P.~Dusart,
\textit{Estimates of some functions over primes without RH},
arXiv:1002.0442 (2010).
\end{thebibliography}

\end{document}
